% AUTORIDAD DE FUENTE: Exposé IX §§12.6--12.7.2, líneas francesas
% 3674--3707; índices de escaneo 486--488, folios 475--477.

%%%%%%%%%%  scan idx 486  |  folio 475  %%%%%%%%%%

\medskip
\noindent
12.6. Obsérvese que este enunciado demuestra 10.4 en el caso particular
considerado aquí cuando, en 10.4 b), se toma como $\xi$ la polarización
canónica considerada allí. Como hemos mostrado en 10.5, esto basta para
deducir 10.4 en el caso general.

\medskip
\noindent
12.7. \underline{Demostración de 12.5. El caso $\ell\ne p$.} Se apoya
esencialmente en la \underline{fórmula cohomológica de
Picard--Lefschetz} en la situación considerada en 12.5, fórmula que da la
acción del grupo de monodromía $I$ sobre el grupo siguiente:

\begin{equation}\tag{12.7.1}
 H^1(X_{\bar\eta},\mathbb Z_\ell)
 \simeq H^1(A_{\bar\eta},\mathbb Z_\ell)
 \simeq\operatorname{Hom}
   \bigl(T_\ell(A_{\bar\eta}),\mathbb Z_\ell\bigr)
\end{equation}

\noindent
cuando $\ell\ne p$. Esta fórmula, que se demostrará en un exposé
posterior (Exposé XV), y de hecho en un caso considerablemente más
general, es esencialmente idéntica a la aserción 12.5 para el primo
$\ell$ en cuestión, teniendo en cuenta la definición de $u_\ell$ dada en
9.2 para $\ell\ne p$. Queda, pues, por tratar el caso $\ell=p$. Lo
haremos mediante un método que lo reduce a la característica cero (y,
además, al caso en que $X_o$ tiene un único punto singular), método que
explicaremos a continuación.

%%%%%%%%%%  scan idx 487  |  folio 476  %%%%%%%%%%

\medskip
\noindent
\underline{Observaciones} 12.7.2. Obsérvese asimismo que, una vez hecha
la reducción anunciada, los argumentos habituales permitirían reducirse
al caso en que el cuerpo base es el cuerpo $\mathbb C$ de los números
complejos e invocar entonces la fórmula de Picard--Lefschetz en su forma
clásica, salvo que no parece existir una referencia satisfactoria para
esta fórmula. La demostración de la fórmula de Picard--Lefschetz que se
dará en nuestro Seminario procederá de manera análoga, por reducción al
caso trascendente, susceptible de métodos propiamente trascendentes.

%%%%%%%%%%  scan idx 488  |  folio 477  %%%%%%%%%%

Puede decirse, por tanto, que, en cualquier caso, la demostración de 12.5
(y, por consiguiente, de 10.4) que proponemos aquí es de naturaleza
trascendente. Señalemos también que la única otra demostración conocida
de 10.4 procede de la teoría analítica rígida [24] y, por esta razón,
debe considerarse igualmente como esencialmente «trascendente». Sería
interesante hallar demostraciones puramente algebraicas de estos
resultados (¡aunque no es en absoluto evidente que ello sea posible!).
